\section{系统概述}

\subsection{设计背景}
学术内容真实性识别已不再局限于传统图像篡改检测场景。随着生成式人工智能技术在论文写作、图片生成与学术评审文本中的广泛使用，学术不端风险逐渐呈现出多类型、跨环节与高隐蔽性的特点。对于期刊、科研机构和研究团队而言，仅依靠单一内容类型、单一检测模型或人工抽查方式，已难以满足对学术资源进行系统化风险识别的需求。

当前项目已完成一轮以学术图像检测为核心的系统建设，具备用户注册登录、图像上传、图片提取、图像自动检测、人工审核、结果通知以及管理端统计等基础能力。从代码实现上看，系统已经形成了由用户端、管理端、Django后端、异步任务系统和独立AI推理服务组成的可运行基础框架，为后续迭代提供了较好的工程基础。

然而，现有系统的业务抽象和数据模型仍明显围绕图像检测链路展开，尚不能完整支撑迭代需求中提出的全篇论文检测、Review检测、综合报告统一生成、模型配置治理以及论文/Review专项日志管理等能力。因此，本轮概要设计的核心任务不是简单补充页面或接口，而是在保留现有可用能力的基础上，完成系统从“图像检测平台”向“综合学术AI鉴伪平台”的结构升级。

\subsection{现有系统基础}
结合当前项目代码，现有系统已经具备如下实现基础：
\begin{itemize}[leftmargin=2em]
    \item 在前端层面，已分别实现用户端与管理端两套Web应用，能够支持基本业务操作、历史记录查询、人工审核入口与管理统计展示；
    \item 在后端层面，已形成以Django为核心的业务服务，提供用户认证、资源上传、检测任务管理、人工审核、组织管理、通知与管理端接口；
    \item 在异步处理层面，已集成Celery与Redis，支持检测任务异步执行、批处理调度、任务状态更新与结果回写；
    \item 在AI推理层面，已具备独立的图像推理服务和图像检测流水线，能够执行图像伪造识别、可疑区域生成与结果封装；
    \item 在数据层面，已具备用户、组织、上传文件、检测任务、检测结果、人工审核、日志与通知等核心表结构。
\end{itemize}

上述基础说明，当前系统并非从零开始设计，而是在已有实现上进行面向迭代需求的重构与扩展。因此，概要设计需要同时考虑两个方面：一方面要充分复用已有实现中成熟的模块和机制；另一方面要识别现有设计中图像中心化、模型耦合和实体抽象不足等问题，并提出统一化、可扩展的改造方案。

\subsection{现有系统与迭代目标差异}
从需求文档与现有实现对照来看，当前系统与本轮目标系统之间主要存在以下差异：
\begin{longtable}{|p{3.8cm}|p{4.8cm}|p{5.4cm}|}
\hline
对比项 & 当前系统现状 & 本轮迭代目标 \\ \hline
核心检测对象 & 以学术图像为主 & 统一支持学术图像、全篇论文与Review \\ \hline
业务抽象中心 & 以图像上传和图像结果为中心 & 以统一学术资源、统一检测任务和统一报告为中心 \\ \hline
AI能力组织方式 & 图像检测链路已落地，模型调用与业务耦合较强 & 不同检测对象对应不同检测适配器，模型调用更加配置化与模块化 \\ \hline
结果输出方式 & 以图像检测结果和图像报告为主 & 统一支持图像、论文、Review结果展示与综合报告生成 \\ \hline
人工审核范围 & 以图像审核流程为主 & 需要与更广义的综合鉴伪流程进行衔接 \\ \hline
日志与治理能力 & 已具备操作日志和部分统计能力 & 进一步增加论文日志、Review日志、模型配置和资源治理能力 \\ \hline
\end{longtable}

因此，本轮概要设计的重点不在于重写已有图像检测流程本身，而在于建立能够承载三类学术资源的总体结构，使现有功能能够作为其中一个可复用子链路纳入统一架构。

\subsection{设计目标}
本轮概要设计的总体目标如下：
\begin{itemize}[leftmargin=2em]
    \item 在保留现有学术图像检测能力的前提下，将系统扩展为面向学术图像、全篇论文和Review的统一鉴伪平台；
    \item 对上传、检测、审核、报告、日志与后台治理形成一致的系统结构，避免不同资源类型各自独立演化；
    \item 降低现有实现中图像中心化带来的耦合问题，为新增资源类型与新增检测模型预留稳定扩展点；
    \item 明确前端、后端、异步任务系统、AI推理服务、数据库和文件存储之间的职责边界与协作关系；
    \item 为后续详细设计、接口扩展、数据库调整和模块重构提供统一依据。
\end{itemize}

进一步地，从工程交付角度看，本轮设计还应满足以下目标：
\begin{itemize}[leftmargin=2em]
    \item 能够支撑逐步迭代开发，而不是要求一次性完成全部重构；
    \item 能够兼容已有图像检测代码与数据结构，降低重构风险；
    \item 能够将新增论文检测和Review检测纳入现有任务调度、通知和管理体系中；
    \item 能够支撑课程项目背景下的文档撰写、演示展示和后续验收。
\end{itemize}

\subsection{用户与角色视角}
根据需求文档和当前实现，本系统主要服务于以下三类角色：
\begin{itemize}[leftmargin=2em]
    \item \textbf{普通用户}：主要包括科研人员、学生、编辑和其他资源提交者。其关注点在于上传资源、发起检测、查看结果、下载报告以及申请人工审核；
    \item \textbf{审核参与者}：主要指具有人工审核职责的用户。其关注点在于接收审核任务、查看资源与自动检测结果、提交审核结论及原因说明；
    \item \textbf{管理员}：包括组织管理员和超级管理员。其关注点在于组织和用户治理、审核审批、资源管理、日志查询、统计分析和模型配置。
\end{itemize}

不同角色共享同一业务平台，但其功能入口、数据访问范围和操作边界不同。因此，概要设计中需要明确角色分层，而不能将系统视为单一用户形态的应用。

\subsection{业务范围}
结合需求规格说明书，本轮迭代后系统的业务范围主要包括以下几个部分：
\begin{itemize}[leftmargin=2em]
    \item \textbf{资源接入范围}：系统支持图像文件、论文文件和Review文本三类资源进入统一检测流程；
    \item \textbf{自动检测范围}：系统支持对图像伪造风险、论文AIGC风险和Review模板化/AI生成倾向进行自动分析；
    \item \textbf{结果输出范围}：系统支持检测结果展示、结构化结果查询与综合报告生成；
    \item \textbf{人工复核范围}：系统支持用户发起人工审核请求，并在必要时引入审核参与者和管理员共同完成复核流程；
    \item \textbf{后台治理范围}：系统支持对用户、组织、资源、日志、审核和模型进行统一管理。
\end{itemize}

\subsection{设计原则}
为了保证本轮系统设计既能兼容当前实现、又能支撑后续演进，概要设计遵循以下原则：
\begin{itemize}[leftmargin=2em]
    \item \textbf{模块化原则}：按业务职责划分模块，减少不同模块之间的直接耦合；
    \item \textbf{统一抽象原则}：对不同类型学术资源采用统一任务与结果抽象，不再仅围绕图像对象组织系统；
    \item \textbf{异步处理原则}：长时检测任务通过异步机制执行，避免阻塞用户交互主流程；
    \item \textbf{可扩展原则}：为新增检测对象、检测模型和管理能力预留扩展接口；
    \item \textbf{可追溯原则}：关键业务过程保留任务状态、日志与报告映射关系，保证结果可回查；
    \item \textbf{兼容演进原则}：在设计上兼顾现有系统可复用部分与目标结构升级，避免一次性推翻现有实现。
\end{itemize}

\subsection{系统边界}
从系统边界看，本系统主要由用户端、管理端、后端业务服务、异步任务处理服务、AI推理服务、数据库和文件存储构成。系统对外提供Web访问与业务接口能力，对内通过HTTP调用、消息队列、数据库记录和文件共享完成各组件之间的协作。

系统边界之外的内容主要包括：
\begin{itemize}[leftmargin=2em]
    \item 具体检测模型的训练过程与训练平台；
    \item 第三方学术数据库或外部内容平台的直接接入；
    \item 大规模生产级运维体系中的弹性调度与云原生编排能力；
    \item 不属于本课程项目实现范围的深度审计或合规系统。
\end{itemize}

也就是说，本文档关注的是系统如何组织、调用和管理检测能力，而不是对底层模型训练机制或外部科研生态系统做全面设计。
